\section{Laplace and Fourier methods for convolution kernels}
\label{app:transforms}

The transform methods of this appendix apply \emph{only} to the convolution case
$K(t,u)=\phi(t-u)$ (no/baseline covariates). Their failure for non-convolution
kernels is exactly what makes the excitation-covariate case
(\S\ref{sec:excitation-cov}) hard.

\subsection{The transforms and the convolution theorem}

The (one-sided) Laplace transform of $f$ is
$\Lap\{f\}(z)=\int_0^\infty e^{-zt}f(t)\dd t$, and the Fourier transform is
$\Four\{f\}(\omega)=\hat f(\omega)=\int_{-\infty}^\infty e^{-i\omega t}f(t)\dd t$.
The one property we need is that both turn convolution into multiplication:
\begin{equation}
\Lap\{\phi\conv\xi\}=\Lap\{\phi\}\,\Lap\{\xi\},\qquad
\widehat{\phi\conv\xi}=\hat\phi\,\hat\xi .
\end{equation}

\begin{intuition}
A convolution mixes every past value of $\xi$ into the present through a
\emph{shift-invariant} weight. The transforms re-express $\xi$ in terms of pure
oscillations $e^{i\omega t}$ (or modes $e^{zt}$), and a shift-invariant system acts
on each such mode merely by scaling it. So in the transform domain the tangled
convolution becomes an independent, per-frequency multiplication --- the system is
``diagonalised.''
\end{intuition}

\subsection{Solving the convolution Volterra equation}

Apply $\Lap$ to $\xi=s+\phi\conv\xi$:
$\Lap\{\xi\}=\Lap\{s\}+\Lap\{\phi\}\Lap\{\xi\}$, hence
\begin{equation}
\Lap\{\xi\}(z)=\frac{\Lap\{s\}(z)}{1-\Lap\{\phi\}(z)},
\qquad
\hat\xi(\omega)=R(\omega)\,\hat s(\omega),\quad R(\omega)=\frac{1}{1-\hat\phi(\omega)} .
\label{eq:transfer}
\end{equation}
The multiplier $R$ is the \emph{transfer function}; inverting the transform
recovers $\xi$. $R=1/(1-\hat\phi)$ is the frequency-domain face of the geometric
series $\sum_n\hat\phi^{\,n}$ --- the resolvent of Appendix~\ref{app:volterra}.

\begin{example}[Exponential kernel]
For $\phi(t)=\kappa\theta e^{-\theta t}$, $\Lap\{\phi\}(z)=\kappa\theta/(z+\theta)$,
so
\begin{equation*}
1-\Lap\{\phi\}=\frac{z+(1-\kappa)\theta}{z+\theta},\qquad
R(\omega)=\frac{\theta+i\omega}{(1-\kappa)\theta+i\omega}.
\end{equation*}
This is a one-pole filter; its pole at $-(1-\kappa)\theta$ is precisely the rate of
the ODE in Appendix~\ref{app:ode}, and inverting $\Lap\{\phi\}/(1-\Lap\{\phi\})$
gives the single-exponential resolvent $h(t)=\kappa\theta e^{(\kappa-1)\theta t}$.
Time domain, frequency domain, and state space are three faces of one operator.
\end{example}

In the multivariate case the same computation gives
$\hat\xi(\omega)=(I-\hat\Phi(\omega))^{-1}\hat s(\omega)$, and at $\omega=0$ the
matrix $\hat\Phi(0)=G$ is the branching matrix, recovering the stationary law
$\bar\xi=(I-G)^{-1}\bar\mu$ (Proposition~\ref{prop:stationary}).

\subsection{The spectral (neural-operator) viewpoint, and why it is exact here}

Because \eqref{eq:transfer} is a pointwise multiply in frequency, the solution
operator $s\mapsto\xi$ is \emph{diagonal} under the Fourier transform. A single
linear Fourier-neural-operator layer computes exactly
$\xi=\Four^{-1}\{R\cdot\Four\{s\}\}$, so with $R=1/(1-\hat\phi)$ it represents the
MBPP operator with \emph{no} approximation; conversely, \emph{learning} $R(\omega)$
from input/output pairs is nonparametric identification of the operator (and hence
of $\phi$ via $\hat\phi=1-1/R$).

\subsection{Why transforms fail under excitation covariates}

All of the above hinges on $K(t,u)=\phi(t-u)$. When the excitation is
covariate-modulated, $K(t,u)=\alpha^0 e^{\delta^\top Z(t)}g(t-u)$ depends on $t$
and $u$ \emph{separately}; the convolution theorem no longer applies, there is no
single transfer function, and the operator is \emph{not} diagonalised by Fourier.
One must return to the time domain --- the LTV ODE or Volterra quadrature of
Appendix~\ref{app:volterra}. This is the precise mathematical price of
\S\ref{sec:excitation-cov}.

\begin{intuition}
Transforms diagonalise \emph{shift-invariant} systems. A covariate that changes the
gain over time breaks shift-invariance --- the system treats ``now'' differently
from ``then'' --- so no single set of frequency weights can describe it. The cure is
to stop trying to diagonalise and instead integrate the dynamics forward in time.
\end{intuition}
